% New English translation from checked cote 25, site 5 (PDF page 6).
% AI-translated by GPT-6 Astra (OpenAI, ChatGPT Pro mode).
\scanpage{25/5}
\mspage{\hand{2}}

or again $\underline{L}^{-1}\subset\underline{E}$) and $\unsure{\phi|\emptyset}$ is obtained by taking \var{\inserted{for $\unsure{\phi|\emptyset}$}}{for $\unsure{\phi|\emptyset}$}{insertion}{ega25_25_005-v0001} a section of $\underline{L}^{-1}$ that
vanishes at no point \var{\xout{\illegible{1}}}{}{typewriter strike-out}{ega25_25_005-v0002}, i.e. a trivialization of $\underline{L}^{-1}$
(which in any case exists locally). Observe also that $H_{\hand{\xi}}$ is
precisely $P(\underline{E}/\underline{L}^{-1})$ (canonical isomorphism), giving a third
way to describe $H_{\hand{\xi}}$ (NB $P(\underline{E}/\underline{L}^{-1})$ is indeed canonically embedded in
$P = P(\underline{E})$), which has the additional advantage of showing that $H$ is a projective
bundle over $S$, and a fortiori is smooth over $S$ (though it would have been necessary to say \var{\xout{\unsure{th}}}{}{typewriter strike-out}{ega25_25_005-v0003}
\var{\margin{$\|$}}{\placedmargin{$\|$}}{marginal note placed}{ega25_25_005-v0004} in par. 17 that a projective bundle is smooth ...). It would no doubt
\var{\margin{$\|$}}{\placedmargin{$\|$}}{marginal note placed}{ega25_25_005-v0005} be preferable to start with this one.
\uline{Remarks}. The formation of $H_{\hand{\xi}}$ in terms of $\hand{\xi}$ is compatible with
base change, as is immediately seen., in other words,
we obtain a homomorphism of functors $(\mathrm{Sch})^{\circ}_{/S}\hand{\to}(\mathrm{Ens})$
\[ \overset{\hand{\vee}}{P} \hand{\longrightarrow} \underline{\mathrm{Div}}\,(P/S) \qquad , \]
where the second member denotes the functor of ``relative divisors'' of $P/S$,
whose value at $S'$ (an arbitrary $S$-prescheme) is the set of
closed subschemes of $P_{S'}$ that are transverse complete intersections
of codimension 1 relative to $S'$ (cf par.19). It is easy to show
that this\var{\xout{\illegible{1}}}{}{typewriter strike-out}{ega25_25_005-v0006} homomorphism of functors is a monomorphism, in other
words that $\hand{\xi}$ is determined by $H_{\hand{\xi}}$. (This last fact justifies the terminology
``scheme of hyperplanes'' used \var{\xout{\unsure{a}}}{}{typewriter strike-out}{ega25_25_005-v0007}above). We shall see \var{\xout{\illegible{1}}}{}{typewriter strike-out}{ega25_25_005-v0008} \var{\margin{)}}{\placedmargin{)}}{marginal note placed}{ega25_25_005-v0009}
\var{\xout{\unsure{Chap.V}}}{}{typewriter strike-out}{ega25_25_005-v0010} that the functor $\underline{\mathrm{Div}}$ $(P/S)$ is representable by \var{\xout{\unsure{a}}}{}{typewriter strike-out}{ega25_25_005-v0011}\var{\inserted{the pre}}{the pre}{insertion}{ega25_25_005-v0012}scheme that is the sum
of the $P(\mathrm{Symm}^{k}(\overset{\hand{\vee}}{E}))$, so that $\overset{\hand{\vee}}{P}$ is identified with an open
and closed subprescheme of $\underline{\mathrm{Div}}$ $(P/S)$ ... (NB To tell the truth, the determination of the relative
\var{\margin{$|$}}{\placedmargin{$|$}}{marginal note placed}{ega25_25_005-v0015} divisors \var{\inserted{of $P/S$}}{of $P/S$}{insertion}{ega25_25_005-v0016} could already be carried out with the means at hand, by
\var{\margin{$|$}}{\placedmargin{$|$}}{marginal note placed}{ega25_25_005-v0017} using the results of Chap III, and could be added \var{\xout{\illegible{1}}}{}{typewriter strike-out}{ega25_25_005-v0018} as
\var{\margin{$|$}}{\placedmargin{$|$}}{marginal note placed}{ega25_25_005-v0019} an example in par.19 ...).
\tnote{25-005-1}{The formula retains the French functor abbreviation Ens. The two French accent corrections in À vrai dire and détermination leave the English expressions unchanged; their variant identifiers are documented in the translation omission record.}

\mspage{\hand{5}}
